% Options for packages loaded elsewhere
\PassOptionsToPackage{unicode}{hyperref}
\PassOptionsToPackage{hyphens}{url}
\PassOptionsToPackage{dvipsnames,svgnames,x11names}{xcolor}
%
\documentclass[
  letterpaper,
  DIV=11,
  numbers=noendperiod]{scrartcl}

\usepackage{amsmath,amssymb}
\usepackage{iftex}
\ifPDFTeX
  \usepackage[T1]{fontenc}
  \usepackage[utf8]{inputenc}
  \usepackage{textcomp} % provide euro and other symbols
\else % if luatex or xetex
  \usepackage{unicode-math}
  \defaultfontfeatures{Scale=MatchLowercase}
  \defaultfontfeatures[\rmfamily]{Ligatures=TeX,Scale=1}
\fi
\usepackage{lmodern}
\ifPDFTeX\else  
    % xetex/luatex font selection
\fi
% Use upquote if available, for straight quotes in verbatim environments
\IfFileExists{upquote.sty}{\usepackage{upquote}}{}
\IfFileExists{microtype.sty}{% use microtype if available
  \usepackage[]{microtype}
  \UseMicrotypeSet[protrusion]{basicmath} % disable protrusion for tt fonts
}{}
\makeatletter
\@ifundefined{KOMAClassName}{% if non-KOMA class
  \IfFileExists{parskip.sty}{%
    \usepackage{parskip}
  }{% else
    \setlength{\parindent}{0pt}
    \setlength{\parskip}{6pt plus 2pt minus 1pt}}
}{% if KOMA class
  \KOMAoptions{parskip=half}}
\makeatother
\usepackage{xcolor}
\setlength{\emergencystretch}{3em} % prevent overfull lines
\setcounter{secnumdepth}{-\maxdimen} % remove section numbering
% Make \paragraph and \subparagraph free-standing
\makeatletter
\ifx\paragraph\undefined\else
  \let\oldparagraph\paragraph
  \renewcommand{\paragraph}{
    \@ifstar
      \xxxParagraphStar
      \xxxParagraphNoStar
  }
  \newcommand{\xxxParagraphStar}[1]{\oldparagraph*{#1}\mbox{}}
  \newcommand{\xxxParagraphNoStar}[1]{\oldparagraph{#1}\mbox{}}
\fi
\ifx\subparagraph\undefined\else
  \let\oldsubparagraph\subparagraph
  \renewcommand{\subparagraph}{
    \@ifstar
      \xxxSubParagraphStar
      \xxxSubParagraphNoStar
  }
  \newcommand{\xxxSubParagraphStar}[1]{\oldsubparagraph*{#1}\mbox{}}
  \newcommand{\xxxSubParagraphNoStar}[1]{\oldsubparagraph{#1}\mbox{}}
\fi
\makeatother


\providecommand{\tightlist}{%
  \setlength{\itemsep}{0pt}\setlength{\parskip}{0pt}}\usepackage{longtable,booktabs,array}
\usepackage{calc} % for calculating minipage widths
% Correct order of tables after \paragraph or \subparagraph
\usepackage{etoolbox}
\makeatletter
\patchcmd\longtable{\par}{\if@noskipsec\mbox{}\fi\par}{}{}
\makeatother
% Allow footnotes in longtable head/foot
\IfFileExists{footnotehyper.sty}{\usepackage{footnotehyper}}{\usepackage{footnote}}
\makesavenoteenv{longtable}
\usepackage{graphicx}
\makeatletter
\def\maxwidth{\ifdim\Gin@nat@width>\linewidth\linewidth\else\Gin@nat@width\fi}
\def\maxheight{\ifdim\Gin@nat@height>\textheight\textheight\else\Gin@nat@height\fi}
\makeatother
% Scale images if necessary, so that they will not overflow the page
% margins by default, and it is still possible to overwrite the defaults
% using explicit options in \includegraphics[width, height, ...]{}
\setkeys{Gin}{width=\maxwidth,height=\maxheight,keepaspectratio}
% Set default figure placement to htbp
\makeatletter
\def\fps@figure{htbp}
\makeatother

\usepackage{graphicx}
\DeclareGraphicsExtensions{.pdf,.png,.jpg,.jpeg}
\KOMAoption{captions}{tableheading}
\makeatletter
\@ifpackageloaded{caption}{}{\usepackage{caption}}
\AtBeginDocument{%
\ifdefined\contentsname
  \renewcommand*\contentsname{Table of contents}
\else
  \newcommand\contentsname{Table of contents}
\fi
\ifdefined\listfigurename
  \renewcommand*\listfigurename{List of Figures}
\else
  \newcommand\listfigurename{List of Figures}
\fi
\ifdefined\listtablename
  \renewcommand*\listtablename{List of Tables}
\else
  \newcommand\listtablename{List of Tables}
\fi
\ifdefined\figurename
  \renewcommand*\figurename{Figure}
\else
  \newcommand\figurename{Figure}
\fi
\ifdefined\tablename
  \renewcommand*\tablename{Table}
\else
  \newcommand\tablename{Table}
\fi
}
\@ifpackageloaded{float}{}{\usepackage{float}}
\floatstyle{ruled}
\@ifundefined{c@chapter}{\newfloat{codelisting}{h}{lop}}{\newfloat{codelisting}{h}{lop}[chapter]}
\floatname{codelisting}{Listing}
\newcommand*\listoflistings{\listof{codelisting}{List of Listings}}
\makeatother
\makeatletter
\makeatother
\makeatletter
\@ifpackageloaded{caption}{}{\usepackage{caption}}
\@ifpackageloaded{subcaption}{}{\usepackage{subcaption}}
\makeatother

\ifLuaTeX
  \usepackage{selnolig}  % disable illegal ligatures
\fi
\usepackage{bookmark}

\IfFileExists{xurl.sty}{\usepackage{xurl}}{} % add URL line breaks if available
\urlstyle{same} % disable monospaced font for URLs
\hypersetup{
  pdftitle={SVD Interlude},
  colorlinks=true,
  linkcolor={blue},
  filecolor={Maroon},
  citecolor={Blue},
  urlcolor={Blue},
  pdfcreator={LaTeX via pandoc}}


\title{SVD Interlude}
\author{}
\date{}

\begin{document}
\maketitle


\section{\texorpdfstring{\emph{A Geometric Buildup to the
SVD}}{A Geometric Buildup to the SVD}}\label{a-geometric-buildup-to-the-svd}

\subsection{\texorpdfstring{\emph{From Input/Output Spaces to Orthogonal
Directions That Stay
Orthogonal}}{From Input/Output Spaces to Orthogonal Directions That Stay Orthogonal}}\label{from-inputoutput-spaces-to-orthogonal-directions-that-stay-orthogonal}

\subsubsection{\texorpdfstring{\emph{Big
Idea}}{Big Idea}}\label{big-idea}

Often a matrix represents a process that takes in one kind of data and
produces another.

Examples:

\begin{itemize}
\tightlist
\item
  3D position \(\to\) 2D screen coordinates
\end{itemize}

. . .

\begin{itemize}
\tightlist
\item
  RGB color (3 numbers) \(\to\) grayscale brightness (1 number)
\end{itemize}

. . .

\begin{itemize}
\tightlist
\item
  3 forces on an object \(\to\) net torque (3 \(\to\) 1)
\end{itemize}

. . .

\begin{itemize}
\tightlist
\item
  5 features \(\to\) a prediction score
\end{itemize}

So a matrix can represent a transformation

\[
A: \mathbb{R}^n \to \mathbb{R}^m
\]

The input space and output space may be \textbf{different worlds}.

. . .

\subsection{\texorpdfstring{\emph{A Concrete Example (3 →
2)}}{A Concrete Example (3 → 2)}}\label{a-concrete-example-3-2}

Consider

\[
A =
\begin{bmatrix}
1 & 0 & 0\\
0 & 1 & 0
\end{bmatrix}.
\]

Then

\[
A(x,y,z) = (x,y).
\]

. . .

Geometrically:

\begin{itemize}
\tightlist
\item
  the component of the input in the \(x\) direction shows up in the
  output
\item
  the component of the input in the \(y\) direction shows up in the
  output
\item
  the component of the input in the \(z\) direction is completely
  ignored
\end{itemize}

. . .

So:

\begin{itemize}
\tightlist
\item
  the entire \(xy\)-plane survives
\item
  everything in the \(z\) direction gets squashed to zero
\end{itemize}

Already we see:

\begin{quote}
Some directions in the input space matter, and some directions don't.
\end{quote}

. . .

\subsection{\texorpdfstring{\emph{Decomposing Any
Vector}}{Decomposing Any Vector}}\label{decomposing-any-vector}

Every vector \(\mathbf{x} \in \mathbb{R}^3\) can be written as

\[
\mathbf{x} = (\text{part in the }xy\text{-plane}) + (\text{part in the }z\text{-direction}).
\]

Apply \(A\):

\[
A\mathbf{x} = A(\text{xy part}) + A(\text{z part}).
\]

But the second term is always zero, so

\[
A\mathbf{x} = A(\text{xy part}).
\]

\begin{quote}
Before doing anything interesting, a transformation throws away the part
of the input it can't see.
\end{quote}

. . .

\subsection{\texorpdfstring{\emph{Example: Rank
1}}{Example: Rank 1}}\label{example-rank-1}

Sometimes a transformation effectively responds to only one input
direction.

For instance, let

\[
A =
\begin{bmatrix}
1 & 1 \\
0 & 0 \\
0 & 0
\end{bmatrix}
: \mathbb{R}^2 \to \mathbb{R}^3.
\]

\subsubsection{\texorpdfstring{\emph{What It
Does}}{What It Does}}\label{what-it-does}

It takes \((x,y)\) and returns

\[
A(x,y) = (x+y,0,0).
\]

\subsubsection{\texorpdfstring{\emph{Multiply It
Out}}{Multiply It Out}}\label{multiply-it-out}

Multiply \(A\) by a vector:

\[
A
\begin{bmatrix}
x \\
y
\end{bmatrix}
=
\begin{bmatrix}
1 & 1 \\
0 & 0 \\
0 & 0
\end{bmatrix}
\begin{bmatrix}
x \\
y
\end{bmatrix}
=
\begin{bmatrix}
x + y \\
0 \\
0
\end{bmatrix}
\]

. . .

Notice that the output depends only on the combination

\[
x+y.
\]

So the component of the input in the direction

\[
\begin{bmatrix}1\\1\end{bmatrix}
\]

affects the output, but the component of the input in the orthogonal
direction

\[
\begin{bmatrix}1\\-1\end{bmatrix}
\]

is completely ignored.

\subsection{\texorpdfstring{\emph{}}{}}\label{section}

So even though the input space is 2-dimensional, the transformation
really only responds to the component of the input in \textbf{one
direction}.

Now:

\begin{itemize}
\tightlist
\item
  all outputs lie along a \textbf{line} in \(\mathbb{R}^3\)
\item
  the map only uses \textbf{one} independent input direction
\end{itemize}

So

\[
\mathrm{rank}(A) = 1.
\]

Geometrically:

\begin{itemize}
\tightlist
\item
  the transformation responds to the component of the input in only one
  direction
\item
  all outputs are constrained to a one-dimensional slice of the output
  space.
\end{itemize}

This will show up later as having only \textbf{one nonzero singular
value}.

. . .

\subsection{\texorpdfstring{\emph{Takeaway}}{Takeaway}}\label{takeaway}

The number of independent directions the map can actually do anything
with is its \textbf{rank} --- regardless of input or output dimension.

\subsection{\texorpdfstring{\emph{From Rank to
SVD}}{From Rank to SVD}}\label{from-rank-to-svd}

Rank tells us \emph{how many} directions a linear map can act on. SVD
tells us \emph{which} directions, and \emph{how strongly}.

\subsection{\texorpdfstring{\emph{The SVD
Motivation}}{The SVD Motivation}}\label{the-svd-motivation}

Let \(r = \mathrm{rank}(A)\). We'll try to pick orthonormal input
directions \(\mathbf{v}_1,\dots,\mathbf{v}_r\) that capture everything
\(A\) can actually do.

Ask:

\begin{quote}
Can we choose input directions \(\mathbf{v}_1,\dots,\mathbf{v}_r\) so
that their images \([A\mathbf{v}_1,\dots,A\mathbf{v}_r]\) are
\textbf{orthogonal} in the output space?
\end{quote}

That would mean the transformation treats these directions
\textbf{independently}.

Define

\[
\mathbf{w}_j = A\mathbf{v}_j.
\]

We want:

\[
\mathbf{w}_i^T \mathbf{w}_j = 0 \quad (i \ne j).
\]

But

\[
\mathbf{w}_i^T \mathbf{w}_j
= (A\mathbf{v}_i)^T(A\mathbf{v}_j)
= \mathbf{v}_i^T A^T A \mathbf{v}_j.
\]

\subsection{\texorpdfstring{\emph{}}{}}\label{section-1}

\[
\mathbf{w}_i^T \mathbf{w}_j
= (A\mathbf{v}_i)^T(A\mathbf{v}_j)
= \mathbf{v}_i^T A^T A \mathbf{v}_j.
\]

The output vectors will be orthogonal if we choose the \(\mathbf{v}_j\)
to be eigenvectors of \(A^T A\).

Let

\[
A^T A \mathbf{v}_j = \sigma_j^2 \mathbf{v}_j.
\]

Then

\[
\|\mathbf{w}_j\|^2
= (A\mathbf{v}_j)^T(A\mathbf{v}_j)
= \mathbf{v}_j^T A^T A \mathbf{v}_j
= \sigma_j^2.
\]

So the outputs are orthogonal, with known lengths.

. . .

\subsection{\texorpdfstring{\emph{Normalize the
Outputs}}{Normalize the Outputs}}\label{normalize-the-outputs}

Define

\[
\mathbf{u}_j = \frac{1}{\sigma_j} \mathbf{w}_j = \frac{1}{\sigma_j} A\mathbf{v}_j.
\]

Then

\[
 \mathbf{u}_1,\dots,\mathbf{u}_r \text{ are orthonormal}
\]

and

\[
A\mathbf{v}_j = \sigma_j \mathbf{u}_j.
\]

. . .

\subsection{\texorpdfstring{\emph{Matrix
Form}}{Matrix Form}}\label{matrix-form}

Let

\[
V_r = [\mathbf{v}_1\ \cdots\ \mathbf{v}_r],
\qquad
U_r = [\mathbf{u}_1\ \cdots\ \mathbf{u}_r],
\qquad
\Sigma_r = \mathrm{diag}(\sigma_1,\dots,\sigma_r).
\]

Then the relations \(A\mathbf{v}_j = \sigma_j \mathbf{u}_j\) become

\[
A V_r = U_r \Sigma_r.
\]

. . .

\subsection{\texorpdfstring{\emph{Geometric
Interpretation}}{Geometric Interpretation}}\label{geometric-interpretation}

For any input \(\mathbf{x}\):

\begin{enumerate}
\def\labelenumi{\arabic{enumi}.}
\tightlist
\item
  express \(\mathbf{x}\) in the orthonormal directions \(\mathbf{v}_j\)
\item
  scale each coordinate by \(\sigma_j\)
\item
  rebuild the result in the orthonormal directions \(\mathbf{u}_j\)
\end{enumerate}

So, in the right orthonormal coordinates in the input and output spaces,

\begin{quote}
Multiplication by \(A\) is just independent scalings of the components
along orthogonal directions in the two spaces.
\end{quote}

. . .

\subsection{\texorpdfstring{\emph{Extending to Full
Bases}}{Extending to Full Bases}}\label{extending-to-full-bases}

We've described what happens to the \(r\) independent directions in the
input space.

. . .

But what if the input space has more than \(r\) dimensions? Or the
output space?

. . .

The trick is to extend the orthonormal bases to full bases.

\begin{itemize}
\tightlist
\item
  Extend \(\mathbf{v}_1,\dots,\mathbf{v}_r\) to an orthonormal basis of
  \(\mathbb{R}^n\)
\item
  Extend \(\mathbf{u}_1,\dots,\mathbf{u}_r\) to an orthonormal basis of
  \(\mathbb{R}^m\)
\end{itemize}

\subsection{\texorpdfstring{\emph{}}{}}\label{section-2}

Let's start with the case where the input space has more than \(r\)
dimensions.

. . .

We started with the \(r\) singular vectors
\(\mathbf{v}_1,\dots,\mathbf{v}_r\). These capture everything \(A\) can
actually do.

. . .

The rest of the input space is the null space of \(A\), and our
additional vectors will be a basis for this null space.

. . .

Let's call these additional vectors
\(\mathbf{v}_{r+1},\dots,\mathbf{v}_n\). We know that

\[
A\mathbf{v}_{r+1} = \mathbf{0}, \quad \dots, \quad A\mathbf{v}_n = \mathbf{0}.
\]

\subsection{\texorpdfstring{\emph{}}{}}\label{section-3}

Similarly, we can extend the singular vectors
\(\mathbf{u}_1,\dots,\mathbf{u}_r\) to a basis for the output space.

. . .

Let's call these additional vectors
\(\mathbf{u}_{r+1},\dots,\mathbf{u}_m\). They represent directions that
cannot be reached by \(A\mathbf{x}\) for any \(\mathbf{x}\).

. . .

This means that the projection of any output \(A\mathbf{x}\) onto any of
these additional vectors must be zero:

\[
\mathbf{u}_{r+k}^T A\mathbf{x} = 0 \text{ for all } \mathbf{x}.
\]

for \(k = 1, \dots, m - r\).

. . .

And the only way this can be true for all \(\mathbf{x}\) is if

\[
\mathbf{u}_{r+k}^T A = \mathbf{0}
\]

\begin{center}\rule{0.5\linewidth}{0.5pt}\end{center}

\subsection{\texorpdfstring{\emph{The transformation in adapted
coordinates}}{The transformation in adapted coordinates}}\label{the-transformation-in-adapted-coordinates}

We have built orthonormal bases:

For the \textbf{input space}:

\[
V = [\mathbf{v}_1\ \cdots\ \mathbf{v}_r\ \mathbf{v}_{r+1}\ \cdots\ \mathbf{v}_n].
\]

For the \textbf{output space}:

\[
  U =
  \begin{bmatrix}
  \mathbf u_1 & \cdots & \mathbf u_r & \mathbf u_{r+1} & \cdots & \mathbf u_m
  \end{bmatrix}.
  \]

\subsection{\texorpdfstring{\emph{}}{}}\label{section-4}

Split them into blocks:

\[
V = \begin{bmatrix} V_r & V_0 \end{bmatrix},
\qquad
U = \begin{bmatrix} U_r & U_0 \end{bmatrix}.
\]

where

\begin{itemize}
\tightlist
\item
  \(V_r\): directions where something happens
\item
  \(V_0\): null space directions
\item
  \(U_r\): directions actually reached
\item
  \(U_0\): directions never reached
\end{itemize}

\subsection{\texorpdfstring{\emph{}}{}}\label{section-5}

Put \(A\) between the two bases:

Consider the matrix

\[
U^T A V.
\]

\[
\begin{bmatrix}
U_r^T \\
U_0^T
\end{bmatrix}
A
\begin{bmatrix}
V_r & V_0
\end{bmatrix}=\begin{bmatrix}
U_r^T A V_r & U_r^T A V_0 \\
U_0^T A V_r & U_0^T A V_0
\end{bmatrix}
\]

\subsection{\texorpdfstring{\emph{}}{}}\label{section-6}

\[
\begin{bmatrix}
U_r^T A V_r & U_r^T A V_0 \\
U_0^T A V_r & U_0^T A V_0
\end{bmatrix}
\]

Each block has a geometric meaning.

\subsection{\texorpdfstring{\emph{Top-right block: null space
dies}}{Top-right block: null space dies}}\label{top-right-block-null-space-dies}

Because \[
A V_0 = 0,
\]

we get

\[
U_r^T A V_0 = 0.
\]

\begin{center}\rule{0.5\linewidth}{0.5pt}\end{center}

\subsection{\texorpdfstring{\emph{Bottom row: unreachable
outputs}}{Bottom row: unreachable outputs}}\label{bottom-row-unreachable-outputs}

Because \[
U_0^T A = 0,
\]

we get

\[
U_0^T A V_r = 0,
\qquad
U_0^T A V_0 = 0.
\]

\begin{center}\rule{0.5\linewidth}{0.5pt}\end{center}

\subsection{\texorpdfstring{\emph{Only one block
survives}}{Only one block survives}}\label{only-one-block-survives}

So the matrix collapses to

\$\$ U\^{}T A V =======

\begin{bmatrix}
U_r^T A V_r & 0 \
0 & 0
\end{bmatrix}

\textbackslash end\{bmatrix\}. \$\$

Everything outside the rank-\(r\) interaction vanishes.

\begin{center}\rule{0.5\linewidth}{0.5pt}\end{center}

\section{\texorpdfstring{\emph{What is the remaining
block?}}{What is the remaining block?}}\label{what-is-the-remaining-block}

By construction,

\[
A\mathbf v_i = \sigma_i \mathbf u_i,
\qquad i=1,\dots,r.
\]

Therefore

\$\$ U\_r\^{}T A V\_r ===========

\begin{bmatrix}
\sigma_1 & & \
& \ddots & \
& & \sigma_r
\end{bmatrix}

\& \& \sigma\_r \textbackslash end\{bmatrix\}. \$\$

Call this matrix \(\Sigma_r\).

\begin{center}\rule{0.5\linewidth}{0.5pt}\end{center}

\section{\texorpdfstring{\emph{The final
form}}{The final form}}\label{the-final-form}

\[
U^T A V =
\begin{bmatrix}
\Sigma_r & 0 \
0 & 0
\end{bmatrix}
= \Sigma.
\]

Multiplying back gives

\[
\boxed{A = U \Sigma V^T.}
\]

\begin{center}\rule{0.5\linewidth}{0.5pt}\end{center}

\section{\texorpdfstring{\emph{Conceptual
punchline}}{Conceptual punchline}}\label{conceptual-punchline}

After choosing the right input and output coordinates, the
transformation does only one thing:

\begin{itemize}
\tightlist
\item
  ignore some input directions,
\item
  stretch \(r\) orthogonal directions,
\item
  never produce motion in the remaining output directions.
\end{itemize}

The SVD is simply the matrix form of this geometric statement.




\end{document}
